\subsection{Validering}
Förbättringsförslagen från avsnitt~\ref{subsec:sfera-utokat} utvärderades i en andra omgång intervjuer med scenariobeskrivningar för att validera IT-artefaktens faktiska värde.
En sammanställning av scenariobeskrivningar och svar för intervjuerna i avsnitt finns i bilaga 2.
\paragraph{Journey package}
Att komplettera JP med information om störningar samt kontaktuppgifter till avsändare (tågklarerare) uppfattades generellt som positiv förbättring.
Två respondenter var skeptiska till ifall det var nödvändigt med information om avsändare.
\paragraph{Blankettförberedelse}
Detta förbättringsförslag mottogs väl av samtliga respondenter.
Det skulle med enkla medel effektivt frigöra tid från tågklarerarna så att de kunde lägga tid på andra uppgifter då blanketthantering är mödosamt, tyckte någon.
En respondent tyckte det var viktigt att signalbeteckning inte togs med som informationspunkt eftersom detta kunde ha negativ trafiksäkerhetsmässig inverkan.
\paragraph{Kontaktuppgifter}
Kontaktuppgifter till aktuell driftledningscentral samt bandel var intressant för förarna och föreföll mer intressant när det nämndes i kontexten tågkörning i utlandet.
Att sitta och leta information var enligt respondenterna umbärligt även om man visste var man skulle hitta informationen.
\lq{}Bra om det finns en knapptryckning bort\rq{}, uppgav en respondent.
\paragraph{Trafikstörningar}
Att få information från ursprungskällan uppfattas som åtråvärt för samtliga respondenter.
Särskilt intressant \emph{när} händelsen inträffade enligt en respondent.
Någon respondent trodde att det skulle vara bra ut servicesynpunkt då resenärer och kunder skulle kunna få information snabbare istället för att det skulle gå ut via järnvägsföretaget som får det av Trafikverkets trafikinformatörer.
%\paragraph{Context awareness}
\paragraph{Kvittens av JP}
Responsen på denna punkt var generellt god, men några undrade om det var nödvändigt då ett avvisande av tilldelad JP ändå skulle leda till ett GSM-R-telefonsamtal.
En notis att ny JP mottagits nämndes dock som potentiellt nyttigt för att göras uppmärksam på att läget i trafiken kan ha förändrats.
En respondent tyckte dock att det var naturligt för att undvika envägskommunikation, även tågklarerarens \lq{}order\rq{} gäller.
\paragraph{Rapportering av balisinformationsfel}
Denna punkt uppfattades som positiv men många hade invändningar att det inte fick påverka lokförarens uppmärksamhetsförmåga då ett balisinformationsfel ofta leder till en kraftig inbromsning och därför kräver förarens fulla uppmärksamhet.
En respondent var negativt inställt till detta och föredrog att notera berörd information på ett fysiskt pappersark för att inte bli störd.
Förslaget säger inget om interaktionsdesignen, bara vad meddelandet ska bestå av.
Flera uppgav dock att de trodde att rapporteringsfrekvensen skulle öka markant om detta infördes, i synnerhet av balisinformationsfel typ 1 som inte resulterar i inbromsning utan endast talar om att det uppstod ett fel i informationsöverföringen från en viss balis.
\paragraph{Plattformsbehov}
Respondenten som var godstågsförare reserverade sig för denna funktionalitet.
Övriga var positivt inställda och gillade konceptet att jobba proaktivt för att minska onödig kommunikation.
Om informationen kan göras tillgänglig för tågets egenskaper istället för att det kommuniceras till varje tågklarerare på varje driftplats vore det ännu bättre enligt en respondent.
\paragraph{Meddelande om plötsligt stopp}
Detta förslag fick generellt god respons då förarna potentiellt skulle slippa att lägga onödig tid på att stå i telefonkö när felsökning av fordon krävs.
Någon uppgav att det utöver vad som stod i scenariot var bra att driftledningscentralen vet att föraren inte är tillgänglig på tågets GSM-R-telefon.
I sådant fall kunde meddelandet kompletteras med telefonnummer till arbetstelefon enligt en respondent.
En annan uppgav att detta skulle kunna resultera i att kommunikation \emph{sker} när det kanske annars inte skulle göra det.
Detta skulle avse när lokförare väljer att felsöka och åtgärda direkt istället för att först ringa och meddela att tåget står stilla.
\paragraph{Meddelande om signal i stopp}
Samtliga respondenter uppskattade detta förslag.
Någon jämförde det med \lq{}den berömda, telefonstyrda signalen\rq{} vilket avser när tågklareraren får ett samtal, ser att tåget står vi en stoppsignal och därför väljer att lägga körbesked i signalen istället för att ens svara.
En respondent föreslog att tågklareraren skulle kunna bekräfta meddelandet med en prognos på hur lång tid tåget skulle hålla för stopp.
En förare uppgav att kommunikationen skulle öka på grund av detta då vissa förare drar sig för att ringa till driftledningscentralen i detta läge då de är rädda för att störa tågklareraren med den typen av information.
\paragraph{Meddelande om extra tidsbehov vid plattform}
Denna funktion var uppskattad och en förare nämnde att det kanske kunde ske vid orderläsningen redan innan tåget hade avgått.
En förare uppgav att tågen med platsbokning har väldigt god kännedom om var resenärer med rullstol kommer gå av och på.
En annan förare uppskattade funktionen då det skulle möjliggöra ett toalettbesök om ingen annan prioritering trumfade förfrågan.
En annan förare uppgav att det var ett bra sätt att förhindra att förare inte skulle hedra sin körprofil på grund av någon anledning som kräver plattformstid.
\paragraph{Meddelande till SOS}
Denna funktion var uppskattad av respondenterna av vilka flera hade varit med om att polis och räddningstjänst hade problem att lokalisera deras tåg.
En förare insisterade på att denna funktionalitet skulle få sån effekt att det borde gå att mäta hur enkelt som helst och uppskattas mycket av polisen.
En förare var lite mer försiktigt inställd och undrade om inte driftledningscentralen redan hade denna information
\paragraph{Meddelande om uppdaterad tåglängd}
Detta förslag var generellt sett accepterat som bra men en förare reserverade sig och uppgav denne trodde mer på samtalskommunikation i frågor om tågets egenskaper.
De andra förarna var positivt inställda och uppgav att man ibland blev uppringd och fick frågan \lq{}Hur lång är du (ditt tåg) idag?\rq{}.
En respondent var skeptiskt till att tågklarerarna blev klokare av att veta tåglängder för vissa tåg då hen inte visste vilken kunskap eller information de hade tillgängligt om plattformslängder.
Särskilt uppskattat var de av två förare som körde X2000 respektive X31 och det innebär stor skillnad i hur många motorvagnsenheter tåget kör med.
En X2000-förare uppgav att hen ringde till varje driftplats där hen hade uppehåll för att meddela tågklareraren om att hen körde med \lq{}mult\rq{}.
\paragraph{Meddelande om att samtal önskas}
Förslaget mottogs positivt av alla respondenter.
Som tidigare nämnt i ett annat scenario uppgav en respondent att förare inte vill vara till besvär för tågklarerarna och kanske därför avstår samtal.
Särskilt uppskattat när man har något att säga men som inte är akut uppgav en respondent.